\section{Summary and Conclusion}\label{sec:trace-summary}

There are infinitely many distinct warping free beam models in the family of \cref{sec:gensec} that exactly trace solutions to Saint~Venant's continuum problems. 
% A trace \(\mathcal{T}\) of a continuum displacement \(\hat{\bm u}\) is an operation:
% \RepeatEquation{eq:trace-map}
% such that
% \RepeatEquation{eq:trace-lift-linear} 
Under the Saint~Venant equivalence class (\cref{def:sv-equivalence}) one has the results:
\begin{itemize}
    \item Traces are physically distinguished by how they manifest boundary fixities on the beam displacement \(\bm{u}\) and rotation \(\bm{\theta}\) fields.
    \item Any Class~I trace (\cref{def:class-1}) is compatible with Type~1 beam finite elements and represented through the enhanced strain section of \cref{sec:enhan-section} through relaxation of the consistency condition \eqref{eq:enhan-cii}. %via a Petrov-Galerkin formulation \eqref{eq:petrov-section}, which preserves the classical section resultants. 
    However, most Class I traces are not variational and result in an asymmetric tangent stiffness.

    \item In Class~I, the elastic moment-curvature response of a section is \emph{not} trace dependent (\cref{rem:bend-trace-indep}), and therefore is uniquely determined solely by the cross sectional geometry and material properties. 

    \item The presence of trace parameters in the shear rigidity through \eqref{eq:shear-correction} reveals that the elastic shear response \emph{is} trace dependent. 

    \item The geometric \citep{cowper1966shear} and energetic \citep{renton1991generalized} treatments of shear  
     both lead to exact continuum solutions in the present framework, and are included in the Class I family of traces. 
    
    \item Class~II traces (\cref{def:class-ii}) are the subset of Class I traces which support the variational condition \eqref{eq:enhan-cii} and furnish a symmetric tangent stiffness.
    \item Class~III traces (\cref{def:class-iii}) are the subset of Class II traces that produce the classical stress resultants \((N,\ShearForce,T,\bm{M})\) under variational condition \eqref{eq:enhan-cii}.%the symmetric Bubnov-Galerkin implementation \eqref{eq:bubnov-section}.
    \item All Class~III traces are equivalent under the Saint~Venant class %(\cref{def:sv-equivalence})
    , and their trace matrix is given uniquely by \eqref{eq:energetic-trace-blocks}.
\end{itemize}

\begin{itemize}
    \item There is a 36-parameter space of particular solutions to Saint Venant's problem that can be consistently modeled by standard beam elements with only enhancements to the resultant model.
    \item A 16 parameter subset (\cref{tab:trace-consts}) of these solutions will admit a physically meaningful kinematic interpretation in the configuration variables \(\bm{u},\bm{\theta}\).
    \item A  4 parameter subset is exactly variationally compatible with the linear model of Box~\ref{box:basic-shear} when no warping fields are incorporated.
\end{itemize}



Under the energy-conjugate trace, a clamped boundary condition 
$\bm{u}(0) = \bm{\theta}(0) = \mathbf{0}$ does not fix a direct geometric average over the section, but rather constrains the energy-conjugate kinematic measures in a way that may be difficult to realize physically. 
An illuminating geometric derivation of this model was put forth by \cite{serpieri2014equivalence}, which reveals it's interpretation as a long-wavelength approximation to the continuum problem of a tip-loaded cantilever that is \emph{fully} restrained at it's root \cite{ladeveze1998new}. 
This problem does not fall under the relaxed Saint~Venant class, which precludes Dirichlet conditions of this nature. 



% 1. are these conclusions theoretically sound
% 2. Am i missing any salient points of practical interest?
% 3. How does one understand the implications of the boundary condition phenomenon on beam element continuity? does trace selection have an impact on compatibility in finite element analysis? eg, if i model a saint-venant transversely loaded cantilever with two beam finite finite elements under a given trace, what does the displacement field look like at the element interface?

